\usepackage[14pt]{extsizes}

\usepackage{cmap} % Улучшенный поиск русских слов в полученном pdf-файле
\usepackage[english,russian]{babel} % Языки: английский, русский

\usepackage[T2A]{fontenc} % Поддержка русских букв
\usepackage[utf8]{inputenc} % Кодировка utf8
%\usefont{T2A}{ftm}{m}{sl} % Основная строчка, которая позволяет получить шрифт Times New Roman

\usepackage[left=30mm,right=10mm,top=20mm,bottom=20mm]{geometry}

\usepackage{makecell}
\usepackage{enumitem}
\usepackage{multirow}
\usepackage[para,online,flushleft]{threeparttable}
\usepackage{caption}
% Работа с изображениями и таблицами; переопределение названий по ГОСТу
\captionsetup[table]{singlelinecheck=false, labelsep=endash}
\captionsetup[table]{justification=raggedright,singlelinecheck=off}
\captionsetup{labelsep=endash}
\captionsetup[figure]{name={Рисунок}}
\captionsetup[lstlisting]{justification=raggedright,singlelinecheck=off} % Работа с листингом
\usepackage[justification=centering]{caption} % Настройка подписей float объектов

% \usepackage{graphicx}
% \usepackage{slashbox}
% \usepackage{diagbox} % Диагональное разделение первой ячейки в таблицах

\usepackage{amssymb}
\usepackage{amsmath}
\usepackage{float}
\usepackage{csvsimple}
\usepackage{enumitem} 
\setenumerate[0]{label=\arabic*)} % Изменение вида нумерации списков
\renewcommand{\labelitemi}{---}

% Переопределение стандартных \section, \subsection, \subsubsection по ГОСТу;
% Переопределение их отступов до и после для 1.5 интервала во всем документе
% ГОСТ section
\usepackage{titlesec}


\titleformat{\section}[block]
{\bfseries\normalsize\filcenter}{\thesection}{1em}{}

\titleformat{\subsection}[hang]
{\bfseries\normalsize}{\thesubsection}{1em}{}
\titlespacing\subsection{\parindent}{\parskip}{\parskip}

\titleformat{\subsubsection}[hang]
{\bfseries\normalsize}{\thesubsubsection}{1em}{}
\titlespacing\subsubsection{\parindent}{\parskip}{\parskip}

% ГОСТ изображения и таблицы
\usepackage{caption}
\captionsetup[figure]{name={Рисунок},labelsep=endash}
\captionsetup[table]{singlelinecheck=false, labelsep=endash}

% Список литературы
\makeatletter 
\def\@biblabel#1{#1 } % Изменение нумерации списка использованных источников
\makeatother

\usepackage{setspace}
\onehalfspacing % Полуторный интервал

\frenchspacing
\usepackage{indentfirst} % Красная строка после заголовка
\setlength\parindent{1.25cm}
\usepackage{multirow}
\renewcommand{\baselinestretch}{1.5}
\def\arraystretch{1.5}%  1 is the default, change whatever you need
\setlist{nolistsep} % Отсутствие отступов между элементами \enumerate и \itemize

% Цвета для гиперссылок и листингов
\usepackage{xcolor}
\usepackage{color} 

\titlespacing*{\chapter}{0pt}{-40pt}{8pt}
\titleformat{\chapter}{\filcenter\LARGE\bfseries}{\thechapter}{20pt}{\LARGE\bfseries}
\titleformat{\section}{\Large\bfseries}{\thesection}{20pt}{\Large\bfseries}

\makeatletter
\renewcommand*{\l@chapter}[2]{
	\ifnum \c@tocdepth>\m@ne
	\addpenalty{-\@highpenalty}
	\vskip 1em \@plus.2em
	\@dottedtocline{0}{0pt}{1.5em}{\bfseries #1}{\bfseries #2}
	\fi
}
\makeatother

\renewcommand*{\arraystretch}{1} % математические матрицы интервал между строками

% \usepackage{ulem} % Нормальное нижнее подчеркивание
% \usepackage{hhline} % Двойная горизонтальная линия в таблицах
% \usepackage[figure,table]{totalcount} % Подсчет изображений, таблиц
% \usepackage{rotating} % Поворот изображения вместе с названием
% \usepackage{lastpage} % Для подсчета числа страниц

% Дополнительное окружения для подписей
% \usepackage{array}
% \newenvironment{signstabular}[1][1]{
	% 	\renewcommand*{\arraystretch}{#1}
	% 	\tabular
	% }{
	% 	\endtabular
	% }

% \usepackage{pgfplots}
% \usetikzlibrary{datavisualization}
% \usetikzlibrary{datavisualization.formats.functions}

\usepackage[justification=centering]{caption} % Настройка подписей float объектов

\usepackage[unicode,pdftex]{hyperref} % Ссылки в pdf
\hypersetup{hidelinks}

\newcommand{\code}[1]{\texttt{#1}}

\usepackage{pdfpages}
\usepackage{totcount}
\usepackage{hyperref}